\section*{摘要}
\addcontentsline{toc}{section}{摘要}

\textbf{目的：}对未经后续修改的原始 \path{frankenstein_v28.pt} 进行一次可追溯的回顾性评价，并把单体、六个非 3AV 复合物的预设主分析以及 17 靶点的探索性下游结果统一到同一篇中文稿件中。

\textbf{方法：}模型文件由 SHA-256 冻结。单体留出集含 151 个 Rosetta 理论结构、1505 个真实序列位点；751 个单体原始 PDB 全部标记为 Rosetta 理论模型，不存在可核验的“Baker 33 真结构”子集。来源审计发现旧标签中有 62 个 Ser 被错误标为 N-甲基化；在不改变模型权重与预测概率的前提下，将评价真值修正为 261 个阳性与 1244 个阴性。甲基化阈值预设为严格大于 0.6。复合物主分析仅纳入 1SFI、3P8F、3WNE、3ZGC、4K1E 和 4KEL，在温度 0.5 下保留至少一个已保存甲基化标记的序列。结构判定要求全复合物 CA 在一次全局对齐后，复合物全局 CA RMSD 与末链环肽的 best-forward 循环移位 CA RMSD 同时严格低于同一阈值。17 靶点 best85 结果因按天然序列恢复率进行 best-of-N 选择，仅作探索性支持。

\textbf{结果：}单体天然化氨基酸恢复率（RAA）为 16.08\%，完整扩展 token 恢复率为 14.88\%。修正标签后，已知天然序列甲基分类 AUC 为 0.900，在阈值 $>0.6$ 时 Accuracy、Precision、Recall 和 F1 分别为 86.18\%、64.02\%、46.36\% 和 53.78\%；端到端 AUC 为 0.908，对应四项指标为 89.10\%、88.80\%、42.53\% 和 57.51\%。端到端二分类甲基召回为 111/261（42.53\%），但同时恢复正确天然氨基酸身份和甲基状态的真实甲基残基仅 20/261（7.66\%）。六复合物在温度 0.5 下共有 544 条原始序列，476 条（87.50\%）满足甲基化保留门且全部匹配结构 RMSD。联合 RMSD $<3\,\angstrom$ 为 16/476（3.36\%，Wilson 95\% CI 2.08--5.39\%），$<5\,\angstrom$ 为 101/476（21.22\%，17.78--25.11\%）；以全部 544 条原始序列为分母，端到端产率分别为 2.94\% 和 18.57\%。17 靶点 best85 的肽受体框架 CA RMSD 均值/中位数为 29.93/40.06~\AA，而肽自对齐 CA RMSD 为 2.65/2.53~\AA；这表明内部形状与结合姿态必须分开解释。其平均结构多样性 $1-\mathrm{TM}$ 为 0.660；预测通透性、置信度与自然化 fixed-pose Rosetta 分数仅作为描述性证据。

\textbf{结论：}原始模型在单体甲基化位点排序上具有辨别力，但在阈值 0.6 下召回率中等，且基础序列恢复率较低。六复合物的严格端到端产率显示，主要瓶颈是环肽在同一复合物坐标系中的结合姿态，而非受体/全局对齐本身。当前证据支持把本结果作为模型能力边界与下一轮改进基线，而不支持宣称已实现高成功率结构设计。单体结构 RMSD、显式全原子甲基 RMSD、结合位点恢复、相对天然能量成功率与稳定性尚无可核验结果，未作填补。

\noindent\textbf{关键词：}N-甲基化环肽；ProteinMPNN；序列设计；甲基化分类；RMSD；端到端产率；数据来源审计

\begin{tcolorbox}[colback=SoftAmber,colframe=Amber,title=读者先看这三句话]
  \begin{enumerate}
    \item 正文单体指标使用“去 padding、Ser 来源修正、阈值严格 $>0.6$”口径；用户提供的两张旧终端截图只在附录进行错误复盘。
    \item 六复合物的 3.36\% 与 21.22\% 是通过甲基化保留门后的条件比例；真正从原始 544 条序列出发的产率是 2.94\% 与 18.57\%。
    \item best85 是利用天然序列恢复率选出的 oracle 代表，不是未知天然序列时可直接获得的部署性能。
  \end{enumerate}
\end{tcolorbox}
